\section{Methodology}
\label{sec:methodology}

\subsection{Notation}

Let $\mathcal{S} = \{s_1, s_2, \ldots, s_T\}$ denote an ordered sequence
of axial CT slices comprising one CT study, where $T$ is the number of
slices and varies across studies.
Each slice $s_t$ is associated with a ground-truth label vector
$\mathbf{y}_t \in \{0,1\}^C$ with $C=6$ hemorrhage classes.
The per-slice score predicted by the Stage-1 CNN ensemble is denoted
$\hat{\mathbf{p}}_t^{(1)} \in [0,1]^C$, and the refined score produced by
the Stage-2 sequence model is $\hat{\mathbf{p}}_t^{(2)} \in [0,1]^C$.
We use $\boldsymbol{\ell}_t^{(1)} \in \mathbb{R}^C$ for the pre-sigmoid
logits from Stage~1.

\subsection{Dataset}

We use the CQ500 dataset~\cite{chilamkurthy2018cq500}, which contains
491 de-identified non-contrast head CT studies collected from emergency
departments in India.
Each study was independently labelled by three board-certified neuroradiologists
across six binary findings: \emph{any intracranial hemorrhage} (used to derive
the \emph{no hemorrhage} class), \emph{epidural}, \emph{intraparenchymal},
\emph{intraventricular}, \emph{subarachnoid}, and \emph{subdural} hemorrhage.
Following the procedure of~\cite{wang2021deeplearning}, a majority vote
(at least 2 out of 3 readers) is used to assign the final study-level label.

Labels are propagated to the slice level: all slices in a study inherit the
study-level multi-hot label vector.
This is an intentional approximation---normal slices in a positive study
receive the same label as pathological ones---which is consistent with the
original challenge formulation~\cite{rsna2019} and simplifies the slice-level
loss computation.
After DICOM extraction, the full dataset yields approximately
\num{171000} labelled slices.
We split studies (not slices) into training, validation, and test sets in a
$70/15/15$ ratio with stratification on \emph{any intracranial hemorrhage},
giving $\approx\num{119800}$ training slices,
$\approx\num{25700}$ validation slices, and
$\approx\num{25700}$ test slices.

\subsection{Preprocessing Pipeline}
\label{sec:preprocessing}

\paragraph{DICOM to Hounsfield Units.}
Raw DICOM pixel values $p$ are converted to Hounsfield Units using the
linear rescaling parameters stored in the DICOM header:

\begin{equation}
    v_{\HU} = p \cdot \mathrm{RescaleSlope} + \mathrm{RescaleIntercept}.
    \label{eq:hu_conversion}
\end{equation}

Values are subsequently clipped to $[-1000,\, 3000]\,\HU$ to discard
out-of-body and implant artefacts.

\paragraph{CT Windowing.}
Three diagnostic windows are applied to map the clipped HU values onto
three independent $[0,\,1]$ channels:

\begin{itemize}
    \item \textbf{Brain window}: centre $= 40\,\HU$, width $= 80\,\HU$
          (optimised for grey/white matter contrast).
    \item \textbf{Subdural window}: centre $= 75\,\HU$, width $= 215\,\HU$
          (optimised for subdural collections).
    \item \textbf{Bone window}: centre $= 600\,\HU$, width $= 2800\,\HU$
          (optimised for skull fractures and calcifications).
\end{itemize}

For window centre $w_c$ and width $w_w$, the normalised value of a
pixel with Hounsfield Unit $v$ is

\begin{equation}
    f(v;\,w_c,\,w_w) =
        \mathrm{clip}\!\left(\frac{v - w_c + w_w/2}{w_w},\;0,\;1\right).
    \label{eq:windowing}
\end{equation}

The three resulting single-channel maps are stacked into a
$3 \times H \times W$ tensor, enabling ImageNet-pretrained RGB networks
to process CT data without architectural modification.
This three-window strategy was first proposed by
Sage and Badura~\cite{sage2020doublebranch} and subsequently adopted
by Wang et al.~\cite{wang2021deeplearning}.
All images are resized to $384 \times 384$ pixels using bilinear
interpolation and channel-wise normalised with ImageNet statistics
($\mu = [0.485, 0.456, 0.406]$,
$\sigma = [0.229, 0.224, 0.225]$).

\paragraph{Remark on Adjacent-Slice Context.}
An alternative 3-channel strategy stacks three consecutive slices
$(s_{t-1}, s_t, s_{t+1})$ as colour channels (\emph{adjacent-slice
stacking}), used as Branch~\#2 in~\cite{sage2020doublebranch}.
Following the recommendation of Wang et al.~\cite{wang2021deeplearning},
we adopt the three-window approach in Stage~1 because it encodes richer
per-slice tissue contrast, while inter-slice spatial context is already
captured more effectively by the BiGRU sequence model in Stage~2.

\subsection{Stage 1: Per-Slice CNN Ensemble}
\label{sec:stage1}

\subsubsection{Backbones}

Three architectures form the ensemble:

\begin{enumerate}
    \item \textbf{DenseNet-121}~\cite{huang2017densely}: dense block
          feature maps fed into an adaptive average-pooling head followed
          by a single linear classifier.
          Feature dimension: 1024; final layer: $\mathrm{Linear}(1024, 6)$.
    \item \textbf{DenseNet-169}~\cite{huang2017densely}: identical structure
          with a wider final feature map.
          Feature dimension: 1664; final layer: $\mathrm{Linear}(1664, 6)$.
    \item \textbf{SE-ResNeXt-101-32$\times$4d}~\cite{hu2018squeeze}: a
          ResNeXt backbone augmented with squeeze-and-excitation channel
          attention, using \texttt{AdaptiveConcatPool2d} (concatenation of
          global average and global max pooling) in the head.
          Feature dimension: $2 \times 2048 = 4096$; classification head:
          $\mathrm{BN} \to \mathrm{Dropout}(0.5) \to \mathrm{Linear}(4096, 512)
          \to \mathrm{ReLU} \to \mathrm{BN} \to \mathrm{Dropout}(0.5) \to
          \mathrm{Linear}(512, 6)$.
\end{enumerate}

Each backbone is pretrained on ImageNet-1K.
The three-channel adjacent-slice input is fed to all backbones without
any architectural change to the first convolutional layer (the pretrained
weights are inherited directly).

\subsubsection{Ensemble Averaging}

Given backbone logits $\boldsymbol{\ell}_t^{(k)}$ for model $k \in \{1,2,3\}$,
the ensemble prediction is the arithmetic mean over sigmoid-transformed logits:

\begin{equation}
    \hat{\mathbf{p}}_t^{(1)} =
        \frac{1}{K} \sum_{k=1}^{K} \sigma\!\left(\boldsymbol{\ell}_t^{(k)}\right),
    \quad K = 3,
    \label{eq:ensemble}
\end{equation}

where $\sigma(\cdot)$ denotes the element-wise sigmoid function.
The logit-level ensemble average is equivalently computed as
$\boldsymbol{\ell}_t^{(1)} = \frac{1}{K} \sum_k \boldsymbol{\ell}_t^{(k)}$
and passed as input to Stage~2.

\subsection{Stage 2: Bidirectional GRU Sequence Model}
\label{sec:stage2}

\begin{figure}[t]
    \centering
    \includegraphics[width=\linewidth]{../../reports/figures/gradcam_iph.png}
    \caption{Grad-CAM activation maps for the \emph{intraparenchymal
             hemorrhage} class.  Stage-1 activations concentrate on the
             hyperdense intraparenchymal region; Stage-2 context propagation
             further sharpens the decision boundary at adjacent slices.}
    \label{fig:gradcam_iph}
\end{figure}

The sequence model operates on the logit sequences
$\{\boldsymbol{\ell}_t^{(1)}\}_{t=1}^{T}$ output by the Stage-1 ensemble for all
slices of one CT study, and optionally on the slice thickness
$\delta_t \in \mathbb{R}$ extracted from the DICOM \texttt{SliceThickness}
tag.
Two parallel sub-networks are combined.

\subsubsection{Sequence sub-network 1 (SM1): FC + BiGRU}

A shared linear projection compresses each per-slice logit vector:
\begin{align}
    \mathbf{h}_t^{\mathrm{fc}} &= \mathrm{ReLU}\!\left(
        W_2\, \mathrm{BN}\!\left(W_1 \boldsymbol{\ell}_t^{(1)}\right)\right)
    \in \mathbb{R}^{32}.
    \label{eq:sm1_fc}
\end{align}

The sequence $\{\mathbf{h}_t^{\mathrm{fc}}\}$ is then processed by a
two-layer bidirectional GRU with hidden size $d = 96$:
\begin{align}
    \overrightarrow{\mathbf{g}}_t,\, \overleftarrow{\mathbf{g}}_t
        &= \mathrm{BiGRU}\!\left(\mathbf{h}_t^{\mathrm{fc}}\right),
    \label{eq:sm1_gru}
\end{align}

and the concatenated hidden state
$[\overrightarrow{\mathbf{g}}_t;\, \overleftarrow{\mathbf{g}}_t] \in \mathbb{R}^{2d}$
is projected to per-slice logits by a linear layer:
$\boldsymbol{\ell}_t^{\mathrm{SM1}} = W_3 [\overrightarrow{\mathbf{g}}_t;\, \overleftarrow{\mathbf{g}}_t]
\in \mathbb{R}^C$.

\subsubsection{Sequence sub-network 2 (SM2): Conv1D + BiGRU with skip}

SM2 concatenates the raw logits with the sigmoid predictions from SM1 and,
when available, the normalised slice thickness:

\begin{equation}
    \mathbf{u}_t = \left[\boldsymbol{\ell}_t^{(1)};\;
                         \sigma\!\left(\boldsymbol{\ell}_t^{\mathrm{SM1}}\right);\;
                         \delta_t / \delta_{\max}\right]
    \in \mathbb{R}^{2C+1},
    \label{eq:sm2_input}
\end{equation}

where $\delta_{\max} = 10\,\mathrm{mm}$ is a normalisation constant.
The sequence $\{\mathbf{u}_t\}$ is transposed to channel-first format
$U \in \mathbb{R}^{(2C+1) \times T}$ and processed by a stack of dilated
one-dimensional convolutions:
\begin{align}
    F_1 &= \mathrm{Conv1d}_{k=5,\,d=1}(U), \nonumber \\
    F_2 &= F_1 + \mathrm{Conv1d}_{k=3,\,d=4}\!\bigl(
               \mathrm{Conv1d}_{k=3,\,d=2}(F_1)\bigr),
    \label{eq:sm2_convs}
\end{align}

where $k$ and $d$ denote kernel size and dilation respectively.
$F_2$ is fed into a second BiGRU, whose output is summed with a pointwise
projection of $F_2$ (skip connection), and projected to per-slice logits
$\boldsymbol{\ell}_t^{\mathrm{SM2}} \in \mathbb{R}^C$.

The final Stage-2 output is the element-wise sum of both sub-networks:
\begin{equation}
    \boldsymbol{\ell}_t^{(2)} =
        \boldsymbol{\ell}_t^{\mathrm{SM1}} + \boldsymbol{\ell}_t^{\mathrm{SM2}}.
    \label{eq:stage2_out}
\end{equation}

\subsection{Loss Function and Class Imbalance}
\label{sec:loss}

Multi-label classification is framed as $C$ independent binary problems.
We use binary cross-entropy with logits, summed over classes and averaged
over the mini-batch:

\begin{align}
    \BCE = -\frac{1}{BC} \sum_{b=1}^{B} \sum_{c=1}^{C}
        &\Bigl[ w_c \, y_{b,c} \log \hat{p}_{b,c} \nonumber \\
        &\quad + (1 - y_{b,c}) \log (1 - \hat{p}_{b,c}) \Bigr],
    \label{eq:bce}
\end{align}

where $B$ is the batch size, $y_{b,c} \in \{0,1\}$ is the ground-truth
label, $\hat{p}_{b,c} = \sigma(\ell_{b,c})$ is the predicted probability,
and $w_c \geq 1$ is the positive class weight for class $c$, computed as
the ratio of negative to positive examples in the training set.
This formulation is equivalent to PyTorch's
\texttt{BCEWithLogitsLoss(pos\_weight=...)} and directly addresses the
severe class imbalance present in CQ500 (e.g.,
epidural hemorrhage prevalence $< 5\,\%$).

\subsection{Training Protocol}
\label{sec:training}

All Stage-1 backbone networks are fine-tuned end-to-end from ImageNet
pretrained weights.
We adopt the training protocol of~\cite{wang2021deeplearning}:
\begin{itemize}
    \item \textbf{Optimiser}: Adam~\cite{kingma2015adam} with learning rate
          $\eta = 5 \times 10^{-4}$, weight decay $\lambda = 2 \times 10^{-5}$,
          and $\epsilon = 10^{-8}$.
    \item \textbf{Scheduler}: Warm-restart cosine annealing (SGDR)
          ~\cite{loshchilov2017sgdr} with initial period $T_{\max} = 5$ epochs
          and $\eta_{\min} = 10^{-5}$.
    \item \textbf{Duration}: 80 epochs with early stopping patience of 20
          epochs, monitoring validation macro-AUC.
    \item \textbf{Input}: Three-channel adjacent-slice images at
          $384 \times 384$ pixels.
    \item \textbf{Batch size}: 16 samples; mixed-precision (FP16) training.
\end{itemize}

Data augmentation during Stage-1 training comprises random horizontal
flipping, random rotation up to $\pm 15^{\circ}$, grid-shuffle, and
brightness/contrast jitter, following the protocol of the original
solution.
Stage-2 is trained for 40 epochs with Adam ($\eta = 10^{-3}$) and cosine
annealing, using the frozen Stage-1 logits extracted over the training set.
